\section{Replication: Haggag and Paci (2014) and the Origins of the Base}
\label{app:replication}

\textcite{haggag2014default} established that default menus cause tips
using New York taxi data from 2009--2013, exploiting the fact that the
city's two payment vendors displayed different suggestions. We replicate
that vendor comparison on six archival months (April--June and October
2012, April--May 2013; 42.7 million credit-card trips). The two vendors
differed not only in the \emph{level} of their suggestions but in the
\emph{base} they computed them on---a fact documented independently by
\textcite{thakral2026five}, whose description of the two bases our match
rates confirm. What we add here is the arithmetic linking that base
difference to the tip gap \textcite{haggag2014default} measured.

\subsection{Recovering the menus}
Figure~\ref{fig:hp} plots the distribution of tips as a percentage of the
metered fare, by vendor, for fares of \$15 and above. Vendor 2 (VTS)
shows sharp spikes at exactly 20.0 and 25.0 percent. Vendor 1 (CMT)
shows spikes displaced slightly to the right, at 20.5 and 25.5 percent.
The displacement is not noise, and it is not a different menu level: it
is what a 20 percent button produces when the percentage is applied to a
base that includes the tax and tolls sitting on top of the fare. Vendor
2's spikes land on round percentages of the fare because its base is the
fare plus the time-of-day surcharge, and that surcharge is zero on the
majority of daytime trips---so for most of the distribution its base and
the metered fare coincide. This is exactly why the next subsection tests
the candidate bases directly rather than reading them off spike
locations: the spikes alone cannot separate ``fare'' from ``fare plus a
surcharge that is usually zero.''

\begin{figure}[!htbp]
  \centering
  \singlespacing
  \caption{\capttitle{Vendor Default Menus in the Archival Data, 2012--2013}}
  \label{fig:hp}
  \includegraphics[width=0.84\textwidth]{fig-hp.pdf}
  \floatnotes{Distribution of credit-card tips as a percentage of the
  metered fare, in half-point bins, for trips with fares of \$15 or more.
  Spikes mark each vendor's on-screen default buttons. Vendor 2's spikes
  fall exactly on 20 and 25 percent of the fare; Vendor 1's are displaced
  upward because its percentages are computed on a surcharge-inclusive
  base.}
\end{figure}

\subsection{Testing the base directly}
Table~\ref{tab:hp-base} makes the point without reference to spike
locations. For each vendor we compute the share of positive card tips
that equal a 20, 25, or 30 percent button applied to each candidate
base, walking up the fare stack one component at a time. Testing the
full ladder matters: the surcharge field is zero on most daytime trips,
so ``fare'' and ``fare plus surcharge'' coincide there, and a test that
compares only the fare against the all-in total will appear to match the
fare for a vendor that in fact uses fare-plus-surcharge.

The pattern is categorical and identical in all six months. Vendor 2
(VTS) matches \emph{fare plus surcharge}---60.7 percent, against 30.4
percent for the fare alone and 1.4 percent for the all-in total. Vendor
1 (CMT) matches the \emph{all-in total including tax and
tolls}---53.8 percent, against essentially nothing for either narrower
base. So the vendors did not disagree about whether to include the
time-of-day surcharge; both included it. They disagreed about the sales
tax and tolls sitting above it. This is the base description in
\textcite{thakral2026five}, and our match rates confirm it.

\begin{table}[!htbp]
  \centering
  \singlespacing
  \caption{\capttitle{Which Base Did Each Vendor Use in 2012--2013?}}
  \label{tab:hp-base}
  \small
  \begin{tabular}{@{}lrcccc@{}}
    \toprule
    & & \multicolumn{4}{c}{Share of tips matching a 20/25/30\% button on:} \\
    \cmidrule(lr){3-6}
    Vendor & Card trips & Fare & $+$ surcharge & $+$ tax & $+$ tolls \\
    \midrule
    \input output/tables/tab-hp-base.tex
    \bottomrule
  \end{tabular}
  \floatnotes{Six archival months (April--June and October 2012,
  April--May 2013), credit-card trips with metered fares of \$4--60 and
  positive tips. A match requires equality to within half a cent. Each
  column adds one component of the fare stack to the base of the
  preceding column. Vendor 2 computes on fare plus surcharge; Vendor 1
  computes on the all-in total.}
\end{table}

\subsection{Does the base difference explain the tip difference?}
Vendor 1's mean tip rate is 20.6 percent of fare; Vendor 2's is 19.7
percent, an observed gap of $0.94$ percentage points. Conditioning on
month and fare bin, the gap is $0.0089$ (SE $0.0011$), reported in
Table~\ref{tab:hp}. That gap has a prediction attached to it, and the
correct prediction follows from the correct base difference: what Vendor
1 includes and Vendor 2 does not is the sales tax and tolls, which
average 6.51 percent of the metered fare. A rider pressing a 20 percent
button therefore tips $0.20 \times 0.0651 = 1.30$ percentage points of
fare more at Vendor 1 than at Vendor 2.

Not every rider presses a button. Scaling by Vendor 1's exact-default
share of 53.8 percent gives an expected gap of $0.70$ percentage
points, against an observed $0.94$. Base inclusion accounts for roughly
three-quarters of the vendor gap; the residual quarter is the level
difference in the menus themselves and the difference in zero-tipping
(2.0 percent of Vendor 1's trips against 3.6 percent of Vendor 2's).
The vendor difference that \textcite{haggag2014default} used to identify
default effects is therefore substantially---though not
wholly---attributable to the base-inclusion mechanism this paper
studies.

\begin{table}[!htbp]
  \centering
  \singlespacing
  \caption{\capttitle{Vendor Gap in Mean Tip Rate, 2012--2013}}
  \label{tab:hp}
  \small
  \begin{tabular}{@{}lc@{}}
    \toprule
    & Tip as share of fare \\
    \midrule
    \input output/tables/tab-hp.tex
    \bottomrule
  \end{tabular}
  \floatnotes{Six archival months, credit-card trips with fares of
  \$4--60. The difference conditions on month-by-fare-bin fixed effects
  and clusters by fare bin. \sym{***}, \sym{**}, \sym{*}: significance at
  1, 5, 10 percent.}
\end{table}

\subsection{Executing their deposited code}
The authors' replication archive (openICPSR 113897) contains their
cleaned 2009 analysis file and estimation code. Running their Table 2
specification verbatim on their data---the \$15 menu-threshold
discontinuity in CMT trips, third-order fare polynomial on each side,
day-of-week, hour, and borough controls, driver fixed effects, standard
errors clustered by fare bin, fares \$5--25---yields a tip jump at the
threshold of \$0.2757 (SE \$0.0060) on 6.22 million trips, with a mean
tip of \$2.22 just below the cutoff: their headline default effect,
regenerated from their archive on our machines. The lineage of the
present paper is then exact: their experiment moves the \emph{menu} at
a fare threshold and tips jump 12 percent of the base; ours moves the
\emph{base} at a clock threshold and tips move by 80 percent of the
prevailing tip rate. Both are the same object---suggested-gratuity
arithmetic doing fiscal work riders do not undo.

\subsection{Convergence}
The historical comparison closes the paper's argument. In 2012 the
market was split: one vendor stopped its base at the fare plus the
time-of-day surcharge, the other carried it all the way to the all-in
total including tax and tolls. By 2025, as Table~\ref{tab:vendor} shows,
all three vendors compute on the inclusive total. The design choice that
was once a point of disagreement between payment systems has become the
industry standard, and with it the channel through which government fees
flow into gratuities. That convergence is what makes the channel a
market-wide fiscal fact rather than one vendor's quirk, and it is the
feature of the modern setting that the archival literature could not
have observed.
